\scanpage{32}
\deleted{\uline{Pureté}}

\deleted{\uline{Proposition} Soit $X$ un préschéma loc. noeth. \uline{régulier}.}

\deleted{Conditions équivalentes}

\deleted{(i) Pour tt sous-préschéma régulier $Y$ de $X$, et tout sous-préschéma régulier fermé $Z$ de $Y$, partout de codim $d$, on a}
\[
\deleted{\left\{\begin{aligned}
\underline H^{\unsure{i|i'}}_Z(A_Y)&=0 &&\text{pour } i\ne2d\\
\deleted{\underline H^{2d}_Z(A_Y)}&\quad\deleted{\text{loc. isomorphe à }A_Y}
\end{aligned}\right.}
\]

% slanted
\margin{\uline{Attention} \uline{\unsure{permettre} des} \uline{extensions} \uline{étales} de $X$\dots}

\deleted{(ii) Comme \deleted{\unsure{ci-dessus|dessus}}, avec $d=1$}

\deleted{(iii) Comme}

\deleted{(iii) \deleted{Comme \unsure{ci|(i)}, avec} \inserted{\unsure{pr|p.}} $Y,Z$ comme en i, et $z$ pt maximal de $Z$, on a}
\[
\deleted{\underline H^i_Z(\unsure{H|A}_Y)_{\unsure{\dot z|\bar z}}=0\quad\text{si }i\ne2d.}
\]

\deleted{\uline{Dém} $\begin{aligned}(\mathrm i)&\Rightarrow(\mathrm{ii})\\&\Rightarrow(\mathrm{iii})\end{aligned}$ \unsure{trivial|triviale}.}

\deleted{$(\mathrm{ii})\Rightarrow(\mathrm i)$ par transitivité, en notant que pour $d=1$, $\underline H^2_Z(A_Y)$ est loc isomorphe à $A_Z$}

\deleted{$(\mathrm{iii})\Rightarrow(\mathrm{ii})$ On est ramené au cas $X$ \deleted{\unsure{loc|$\subset$}} $Y$ \inserted{strict} (local \deleted{hensélien}, $Z$ \deleted{\unsure{et de}} $=1$) À priori on a grâce à (iii) que $\underline H^i_Z(A_X)=0$ si $i\ne2$. Récurrence}

% archive sheet number: 31
